% ============================================================
\section{Introduction}\label{sec:intro}
% ============================================================

\subsection{The equations}
We study the motion of two incompressible fluids in a two-dimensional, horizontally periodic porous medium (or Hele-Shaw cell) under gravity.
The fluids have different densities and different viscosities. Their evolution is governed by the transport equation for the density
together with Darcy's law, the \emph{incompressible porous media equation with viscosity jump}:
\begin{equation}\label{eq:IPM}
\left\{\begin{aligned}
&\partial_t\rho+\dvg(\rho u)=0,\\
&\dvg u=0,\\
&\mu(\rho)\,u=-\nabla p-\rho\,\ee,
\end{aligned}\right.
\qquad (x,t)\in \T\times\R\times(0,T),
\end{equation}
where $\T=\R/2\pi\mathbb{Z}$, $x=(x_1,x_2)$, $\ee=(0,1)$, $u$ is the velocity, $p$ the pressure, $\rho$ the density and $\mu(\rho)$ the viscosity.
We have normalised the permeability and the gravity constant to one. Two-phase flows correspond to $\rho$ taking only two values.
After subtracting a hydrostatic pressure and rescaling space, time and pressure, we may assume that
\begin{equation}\label{eq:normalisation}
\rho\in\{-1,1\},\qquad \mu(\rho)=1+A\rho,\qquad A=\frac{\mu^+-\mu^-}{\mu^++\mu^-}\in(-1,1),
\end{equation}
where $\mu^{+}=1+A$ is the viscosity of the denser fluid ($\rho=1$) and $\mu^-=1-A$ that of the lighter one ($\rho=-1$); $A$ is the
viscosity Atwood number. The case $A=0$ is the usual IPM equation. Since $\mu$ is affine in $\rho$ and $\rho^2=1$,
the momentum $m=\rho u$ allows to write Darcy's law \emph{linearly} as $u+Am=-\nabla p-\rho\ee$; this observation, which goes back to
\cite{Men22}, is the starting point of the convex integration theory for \eqref{eq:IPM}.

When the density is $\rho=1$ above a curve $\Gamma(t)$ and $\rho=-1$ below it, \eqref{eq:IPM} reduces to a free boundary problem for $\Gamma(t)$,
the \emph{Muskat problem} \cite{Mus34}. Its well-posedness is governed by the sign of the jump of the normal derivative of the pressure across the interface,
\begin{equation}\label{eq:STfunction-intro}
\sig:=\partial_n p^+-\partial_n p^-\qquad\text{on }\Gamma(t),
\end{equation}
where $n$ is the upward unit normal and $p^\pm$ are the traces of the pressure from above and below. We call $\sig$ the \emph{Saffman--Taylor function}.
If $\sig>0$ (the Rayleigh--Taylor, or Saffman--Taylor, condition) the problem is locally well-posed in Sobolev spaces
\cite{Amb04,CG07,CCG11,Mat19,NP20}; see the surveys \cite{Gan17,GBL20}. If $\sig<0$ the problem is ill-posed in the sense of Hadamard
\cite{SCH04,CG07}: this is the Saffman--Taylor instability \cite{ST58,Hom87}, responsible for the fingering patterns observed when a less viscous fluid
displaces a more viscous one, or when a heavier fluid lies on top of a lighter one. In our normalisation (see Proposition~\ref{prop:sigma})
\begin{equation}\label{eq:sigma-UN}
\sig=-2\,(n_2+A\,U),
\end{equation}
where $U=u\cdot n$ is the normal velocity of the interface. The term $-2n_2$ is the gravitational (Rayleigh--Taylor) contribution and the term $-2AU$ the
viscous (Saffman--Taylor) one: an interface is destabilised by gravity where the heavy fluid is on top ($n_2>0$), and by viscosity where the less viscous fluid
pushes into the more viscous one ($AU>0$). We say that Muskat initial data are \emph{fully unstable} if $\sig<0$ along the whole initial interface.

\subsection{Mixing solutions}
In the unstable regime one does not expect classical solutions, and one is led to look for \emph{weak} solutions which describe the
formation of a \emph{mixing zone}, a region where the two fluids mix at small scales. Otto \cite{Ott99} proposed a relaxation of
\eqref{eq:IPM} (for flat interfaces) based on its gradient-flow structure, and showed that the mixing zone grows linearly in time.
Using the convex integration method of De Lellis and Sz\'ekelyhidi \cite{DLS09} (see also \cite{KMS03}), C\'ordoba, Faraco and Gancedo \cite{CFG11} proved
non-uniqueness of weak solutions of IPM (see also \cite{Shv11}), and Sz\'ekelyhidi \cite{Sze12} computed the relaxation of IPM with $A=0$ and constructed
\emph{mixing solutions} for the flat unstable interface, recovering Otto's linear growth. For non-flat interfaces with $A=0$, Castro, C\'ordoba and Faraco \cite{CCF21}
constructed mixing solutions for every fully unstable $H^5$ interface; the mixing zone grows linearly around an evolving \emph{pseudo-interface} whose
evolution is governed by a delicate nonlocal equation, analysed with a semiclassical pseudodifferential calculus (see also \cite{ACF21}).
F\"orster and Sz\'ekelyhidi \cite{FS18} found a much shorter route via \emph{piecewise constant} subsolutions which only involve the velocity and the acceleration of the
interface at the initial time; Noisette and Sz\'ekelyhidi \cite{NS21} allowed variable opening speeds. Further developments include degraded mixing solutions \cite{CFM19},
localised mixing zones for bubbles and turned interfaces in the partially unstable regime \cite{CFM22}, and related results for the Rayleigh--Taylor instability and
for vortex sheets \cite{GKS21,MS23}. All these works concern equal viscosities, $A=0$.

The case of different viscosities was addressed by Mengual \cite{Men22}, who proved an $h$-principle for \eqref{eq:IPM} with $|A|<1$ (weak solutions exist whenever a
\emph{subsolution} exists), computed the relevant relaxation, constructed mixing solutions for the \emph{flat} unstable interface, and observed that the
relaxation has a ``pinch'' singularity which prevents mixing where there is neither Rayleigh--Taylor instability nor vorticity at the interface.

The purpose of this paper is to construct mixing solutions with viscosity jump for \emph{non-flat}, fully unstable Muskat data. Such a result was announced by Luna and Mengual \cite{LM26}; the present text is an independent reconstruction written from the abstract of that announcement (see the note on the first page). The main difficulty with respect to
\cite{CCF21,FS18,NS21} is that, when $A\neq0$, the velocity is no longer an explicit Birkhoff--Rott integral of the density: already for the sharp interface the
vorticity is only implicitly determined by an integral equation (see \eqref{eq:classical} below), and for subsolutions the situation is worse, because the vorticity of $u$
involves the curl of the (unknown) relaxed momentum.

\subsection{Main result}
Let $\Zc^{k,\gamma}$ denote the set of horizontally periodic, non-self-intersecting $C^{k,\gamma}$ curves $z:\R\to\R^2$ with $z(\alpha)-(\alpha,0)$ $2\pi$-periodic
(Definition~\ref{def:curves}); $n_0$ is the upward unit normal of $z_0\in\Zc^{k,\gamma}$ and $\Omega^\pm(z_0)$ the regions above and below $z_0$. The Muskat data are
\begin{equation}\label{eq:Muskat-data}
\rho_0=\mathbf 1_{\Omega^+(z_0)}-\mathbf 1_{\Omega^-(z_0)}.
\end{equation}
In Proposition~\ref{prop:sigma} we show that the Saffman--Taylor function $\sig_0$ and the normal velocity $U_0$ of the Muskat problem at $t=0$ are given by
\begin{equation}\label{eq:sigma0-intro}
\sig_0=-2\big(I-2A\Kop(z_0)\big)^{-1}[n_{0,2}],\qquad U_0=-\Kop(z_0)\sig_0,
\end{equation}
where $\Kop(z_0)$ is the adjoint Neumann--Poincar\'e operator of the curve (Definition~\ref{def:K}). A weak solution of \eqref{eq:IPM} is a pair
$(\rho,u)\in L^\infty(\T\times\R\times(0,T))^{1+2}$ with $\rho^2=1$ a.e.\ such that the three equations in \eqref{eq:IPM} hold in the sense of distributions on
$\T\times\R\times[0,T)$ with initial datum $\rho_0$, where in the third one $\mu(\rho)u=u+A\rho u$ and the pressure is eliminated by taking the curl
(see \eqref{eq:weak}).

\begin{theorem}\label{thm:main}
Let $A\in(-1,1)$, $\gamma\in(0,1)$ and $z_0\in\Zc^{6,\gamma}$. Assume that the Muskat data \eqref{eq:Muskat-data} are fully unstable:
\begin{equation}\label{eq:fully-unstable}
\sig_0(\alpha)<0\qquad\text{for every }\alpha\in\T .
\end{equation}
Fix $\rhom\in(-1,1)$ and a function $\kappa\in C^{4,\gamma}(\T)$ with
\begin{equation}\label{eq:kappa-range}
0<\kappa(\alpha)<\frac{-\sig_0(\alpha)}{2(1-A\rhom)}\qquad\text{for every }\alpha\in\T.
\end{equation}
Then there exist $T>0$ and curves $Z(\cdot,t)\in\Zc^{3,\gamma}$, $t\in[0,T]$, of the form
\begin{equation}\label{eq:central}
Z(\alpha,t)=z_0(\alpha)+\Big(tU_0(\alpha)+\tfrac{t^2}{2}\,a(\alpha)\Big)n_0(\alpha)+h(t)\,\ee,\qquad a\in C^{3,\gamma}(\T),\quad |h(t)|\le Ct^3,
\end{equation}
with the following property. Let
\[
z^\pm(\alpha,t):=Z(\alpha,t)\pm t\,(1\pm\rhom)\,\kappa(\alpha)\,n_0(\alpha),
\]
let $\Gamma^\pm(t)$ be the curves parametrised by $z^\pm(\cdot,t)$ and let $\Omega_{\rm mix}(t)$ be the open strip between them. Then there exist infinitely many weak
solutions $(\rho,u)$ of \eqref{eq:IPM} on $[0,T]$ with initial datum $\rho_0$ such that
\begin{enumerate}[label=\textup{(\roman*)}]
\item $\rho(x,t)=1$ for a.e.\ $x$ above $\Gamma^+(t)$ and $\rho(x,t)=-1$ for a.e.\ $x$ below $\Gamma^-(t)$;
\item $\rho$ \emph{mixes} in $\Omega_{\rm mix}$: for every open set $B\subset\{(x,t):0<t<T,\ x\in\Omega_{\rm mix}(t)\}$ both $\{\rho=1\}\cap B$ and $\{\rho=-1\}\cap B$ have
positive measure;
\item the coarse-grained density in $\Omega_{\rm mix}$ is $\rhom$: there is a sequence of such solutions converging weakly-$*$ in $L^\infty$ to the function equal to $1$ above
$\Gamma^+(t)$, $\rhom$ in $\Omega_{\rm mix}(t)$, and $-1$ below $\Gamma^-(t)$.
\end{enumerate}
\end{theorem}

Thus the mixing zone has width $2t\kappa(\alpha)$ in the direction $n_0(\alpha)$; it grows linearly in time around the evolving central interface $Z(\cdot,t)$,
which starts with the normal velocity $U_0$ of the Muskat problem, and its upper and lower boundaries move away from $Z$ with normal speeds $(1+\rhom)\kappa$ and $(1-\rhom)\kappa$.
By \eqref{eq:kappa-range} the opening rate $2\kappa$ is bounded by $-\sig_0/(1-A\rhom)$, i.e.\ by the Saffman--Taylor function: the more unstable the interface,
the faster the mixing zone may open. Note that the central interface $Z$ is the ``centre of mass'' of the mixing zone, see Remark~\ref{rem:centre-of-mass}.

\begin{remark}
\begin{enumerate}[label=(\alph*),leftmargin=*]
\item For $A=0$ one has $\sig_0=-2n_{0,2}$, so \eqref{eq:fully-unstable} says that $z_0$ is a graph with the heavy fluid on top, and Theorem~\ref{thm:main} reduces to a result
of the type of \cite{CCF21,FS18,NS21} (with more regularity than needed there). For $A\neq0$ condition \eqref{eq:fully-unstable} is not a priori a graph condition, since the viscous
term $-2AU_0$ in \eqref{eq:sigma-UN} may compensate $n_{0,2}$.
\item The admissible range \eqref{eq:kappa-range} is the one allowed by a subsolution with a single intermediate density $\rhom$. In Section~\ref{sec:remarks} we show (formally)
that subsolutions with continuous profiles allow the upper and lower boundaries to move away from the central interface with normal speeds arbitrarily close to
$-\sig_0/(1-A)$ and $-\sig_0/(1+A)$ respectively (with a common factor $\theta<1$); for the flat interface ($\sig_0\equiv-2$) this gives the asymmetric speeds $2/(1-A)$ and $2/(1+A)$ of Example~\ref{ex:flat}.
\item We have not tried to optimise the regularity of $z_0$.
\end{enumerate}
\end{remark}

\subsection{The key idea: the Saffman--Taylor function as unknown}
Let us explain the main new ingredient. In the classical treatment of the Muskat problem with viscosity jump \cite{CCG11} one writes the velocity as the Birkhoff--Rott
integral of its vorticity $\omega=\gamma\,\delta_\Gamma$, $\gamma=\jump{u\cdot\tau}$. Darcy's law and the continuity of the pressure then give the implicit equation
\begin{equation}\label{eq:classical}
\big(I+2A\Kop(z)\big)\gamma=-2\tau_2,
\end{equation}
see Remark~\ref{rem:duality}. We proceed differently and work with the \emph{Darcy field}
\[
v:=\mu(\rho)\,u=-\nabla p-\rho\ee .
\]
Its jump across the interface is $\jump{v}=-\jump{\nabla p}-\jump{\rho}\ee=-\sig n-2\ee$, because the pressure is continuous. Hence $v$ is the Cauchy integral of the density
$-\sig n-2\ee$, which is \emph{linear in the Saffman--Taylor function}, and the continuity of $u\cdot n=v\cdot n/\mu$ becomes the equation
$(I-2A\Kop(z))\sig=-2n_2$ in \eqref{eq:sigma0-intro}. In other words, we invert the Birkhoff--Rott operator in terms of $\sig$ instead of $\gamma$.

The advantage becomes decisive for subsolutions. For a relaxed state $(\rho,u,m)$ the Darcy field $v=u+Am$ satisfies
\[
\curl v=-\partial_1\rho,\qquad \dvg v=-A\,\partial_t\rho,
\]
so $v$ is \emph{explicitly} determined by $\rho$ and $\partial_t\rho$, whereas the vorticity of $u$ is not (it involves $\curl m$). For a piecewise constant subsolution with
two boundary curves $z^\pm$ moving with normal velocities $V^\pm$, the field $v$ is the sum of two Cauchy integrals with densities
$\jump{\rho}^\pm\big(-\ee+(n_2^\pm+AV^\pm)n^\pm\big)$, and the relaxation conditions only involve $v$, the normal velocities and a divergence-free field in the strip.
When the strip collapses, the relevant combination of the densities is precisely $-\sig n-2\ee$, the opening rate is controlled by $\sig$, and the condition that the
centre of mass of the mixing zone moves with the fluid is, at every order of the time expansion, an equation of the form $(I-2A\Kop(z_0))X=Y$.
Since $\Kop(z_0)$ is a smoothing operator, $I-2A\Kop(z_0)$ is a compact perturbation of the identity, and its invertibility follows from a short energy argument.

\subsection{Strategy}
The proof has four steps.
\begin{enumerate}[label=\textbf{Step \arabic*.},leftmargin=*]
\item \emph{Relaxation} (Section~\ref{sec:relax}). In the variables $(\rho,v,m)$ we compute the first and second order laminates of the constitutive set and obtain the open set
\[
\Uc_A=\big\{(\rho,u,m):\ |\rho|<1,\ |D-a|<|a|\big\},\qquad D=\tfrac{2[(1-A^2)m-(\rho-A)v]}{1-\rho^2}+2\ee,\quad a=\ee+\tfrac{A}{1+A\rho}\,v,
\]
which is contained in the interior of the $\Lambda$-convex hull. Subsolutions with values in $\Uc_A$ yield mixing solutions by the $h$-principle of \cite{Men22}.
\item \emph{Ansatz} (Section~\ref{sec:ansatz}). Given two curves $z^\pm$ bounding a strip with density $\rhom$, we define $v$ as a sum of Cauchy integrals, the relaxed momentum
outside the strip by the constitutive relation, and inside the strip through a stream function interpolating the fluxes prescribed by the Rankine--Hugoniot conditions.
The linear equations then hold \emph{exactly}, for any curves which conserve the total mass. The only remaining condition is the open inequality $|D-a|<|a|$, and it is
controlled by a scalar \emph{flux mismatch} $M(\alpha,t)$.
\item \emph{Expansion in time} (Section~\ref{sec:proof}). Following \cite{FS18}, we take curves which are second order polynomials in $t$. The vanishing of $M$ at order $t^0$ is
the Muskat equation \eqref{eq:sigma0-intro} for the normal velocity; at order $t^1$ it is an equation $(I-2A\Kop(z_0))a=R_1$ for the normal acceleration $a$.
A small vertical translation restores exact mass conservation, and then $M=O(t^2)$.
\item \emph{Conclusion}. With $M=O(t^2)$ the state in the strip is $O(t^\gamma)$-close to its limit as $t\to0$, which lies in $\Uc_A$ exactly when \eqref{eq:kappa-range} holds.
\end{enumerate}

\subsection*{Organisation}
Section~\ref{sec:muskat} introduces periodic Cauchy integrals, the operator $\Kop(z)$ and the Saffman--Taylor formulation of the Muskat problem. Section~\ref{sec:relax}
contains the relaxation and the $h$-principle, together with the explicit flat case. Section~\ref{sec:ansatz} introduces the piecewise constant ansatz and Section~\ref{sec:proof}
proves Theorem~\ref{thm:main}. Section~\ref{sec:remarks} contains remarks on optimal opening rates and open problems. The appendix collects the properties of Cauchy integrals
that we use.
